\documentclass[11pt]{article}

\usepackage[margin=0.85in]{geometry}
\usepackage{booktabs}
\usepackage{tabularx}
\usepackage{array}
\usepackage{amsmath}
\usepackage{graphicx}
\usepackage{xcolor}
\usepackage{enumitem}
\usepackage{microtype}
\usepackage{fancyhdr}
\usepackage{hyperref}
\usepackage[authoryear,round]{natbib}

\hypersetup{
  colorlinks=true,
  linkcolor=blue!55!black,
  urlcolor=blue!60!black,
  citecolor=blue!55!black
}

\pagestyle{fancy}
\fancyhf{}
\lhead{\small Literature Review --- Thermal Extremes \& CVD Admissions}
\rhead{\small Bob Shen / Laidlaw}
\cfoot{\thepage}
\renewcommand{\headrulewidth}{0.3pt}

\setlength{\parindent}{0pt}
\setlength{\parskip}{0.55em}
\setlist[itemize]{leftmargin=1.35em, itemsep=0.2em, topsep=0.2em}
\setlist[enumerate]{leftmargin=1.5em, itemsep=0.2em, topsep=0.2em}

\newcommand{\adds}{\textcolor{green!40!black}{\textbf{Adds.}}}
\newcommand{\inspires}{\textcolor{blue!55!black}{\textbf{Inspires.}}}
\newcommand{\debates}{\textcolor{orange!70!black}{\textbf{Debates.}}}
\newcommand{\impl}{\textcolor{purple!55!black}{\textbf{Implementation.}}}

\title{
  \textbf{Critical Literature Review}\\[0.3em]
  \large What Each Source Adds, Inspires, or Debates for a Monthly Study of\\
  Thermal Extremes and Cardiovascular Admissions in Hong Kong, 2013--2023
}
\author{
  Bob Shen Ruililin\\
  Laidlaw Scholars Programme, The University of Hong Kong\\[0.4em]
  \small Supervisor: Professor David Bishai
}
\date{10 July 2026}

\begin{document}
\maketitle

\begin{center}
\small\textit{
Quality over quantity. This note is not a catalogue of every temperature--health paper.
It is a working map from a small set of high-leverage sources to concrete choices in
research design, modelling, and the project's R pipeline.
}
\end{center}

\tableofcontents
\vspace{0.8em}

%----------------------------------------------------------------------
\section{Purpose and reading guide}
%----------------------------------------------------------------------

The project asks whether \emph{officially defined} thermal extremes---especially hot
nights and cold days---are associated with monthly principal-diagnosis inpatient
admissions for acute myocardial infarction (AMI), ischemic stroke, and hemorrhagic
stroke among Hong Kong residents aged 35+, during 2013--2023, in an aging city that
still experiences cold.

That question sits inside a crowded Hong Kong literature. Claiming ``first'' is
unsafe. Claiming ``nothing is known'' is false. The useful task is sharper:
identify what prior work already settles, what it leaves open, and what it forces
us to change in our own design and code.

Each core source below is read through four lenses:

\begin{itemize}
  \item \adds{} Empirical or conceptual content we should treat as established
        background.
  \item \inspires{} A design idea, metric, or modelling habit worth adopting or
        adapting.
  \item \debates{} A tension with our current plan, a novelty constraint, or a
        claim we must weaken.
  \item \impl{} Concrete implications for \texttt{config.yml}, scripts, data
        requests, or interpretation rules.
\end{itemize}

A recurring intellectual distinction runs through the review:

\begin{quote}
\textbf{Trigger versus burden.} Daily distributed-lag models estimate acute
\emph{triggering} of events over specific lags. Monthly extreme-count models
estimate broader hospital-\emph{burden} associations. Both are scientifically
legitimate; they answer different questions and produce non-comparable
coefficients \citep{bhaskaran2013,gasparrini2010}.
\end{quote}

%----------------------------------------------------------------------
\section{Hong Kong historical baseline: cold-dominant cardiovascular risk}
%----------------------------------------------------------------------

\subsection{Goggins et al.\ (2013) --- AMI, weather, and pollution}

\adds{} In Hong Kong (2000--2009), daily AMI admissions rose as temperatures cooled
below an estimated $\sim$24$^\circ$C threshold ($\sim$3.7\% per 1$^\circ$C using
multi-day lags). A statistically significant heat effect was not the main story in
that decade. Same-day NO$_2$ was the strongest pollutant predictor among those
examined \citep{goggins2013}.

\inspires{} Keep cold as a \emph{co-primary} exposure, not a historical footnote.
Treat heat as an open empirical question rather than a presumed replacement for
cold. Retain NO$_2$ in staged pollution models for AMI.

\debates{} This paper cannot be used as a numeric benchmark for our monthly rate
ratios. Different decade, daily mean temperature, and lag structure. A ``regime
shift from cold to heat'' claim would require methodologically harmonized
comparisons that we do not currently have.

\impl{} Primary model should include cold days alongside hot nights
(\texttt{scripts/11\_*.R} already does this). Do not enter highly collinear heat
metrics together. Avoid direct numeric comparison tables against Goggins
per-$^\circ$C estimates in manuscripts.

\subsection{Goggins et al.\ (2012) --- stroke subtypes separately}

\adds{} Hemorrhagic stroke (HS) showed a strong, approximately linear inverse
association with mean temperature across the observed range (lags 0--4;
$\sim$2.7\% higher admissions per 1$^\circ$C lower). Ischemic stroke (IS) was weaker
and mainly apparent below $\sim$22$^\circ$C with longer lags (0--13). Pollutants were
not associated with either subtype. Effects were stronger among older adults and
women for HS \citep{goggins2012stroke}.

\inspires{} Never pool ``stroke'' as the primary endpoint. Model AMI, IS, and HS
separately. Expect different lag biology even if our monthly design cannot resolve
daily lags: HS may be more acute; IS more delayed/cold-threshold.

\debates{} Monthly aggregation may attenuate short-lag HS signals more than
longer-lag IS/cold signals. A null monthly HS--heat association would not refute
daily HS--cold biology.

\impl{} Keep diagnosis-stratified models mandatory (already planned). Prefer
principal diagnosis. Request sex-specific extracts. Age bands should separate
65--69 and 70--74 rather than collapsing into ``65+''.

\subsection{Lam et al.\ (2018) --- diabetes as effect modifier for AMI}

\adds{} Among public-hospital AMI admissions (2002--2011), patients with diabetes
showed stronger cold-season temperature associations than patients without
diabetes, and some hot-season sensitivity that non-DM patients lacked
\citep{lam2018amiDiabetes}.

\inspires{} Vulnerability is not only chronological age. Comorbidity can reshape
the temperature--AMI curve.

\debates{} Aggregate monthly extracts may not support diabetes stratification.
If comorbidity is unavailable, we must not pretend age alone captures the Lam
finding.

\impl{} Ask Roro whether a diabetes flag (or DM-related secondary diagnosis) is
feasible as a sensitivity stratum. If not, document this as a known limitation in
the assumption ledger (A27) rather than inventing proxies.

%----------------------------------------------------------------------
\section{The decisive recent overlap: hot nights are not a blank slate}
%----------------------------------------------------------------------

\subsection{Guo et al.\ (2024) --- official hot nights versus nighttime heat intensity}

This is the single most important paper for our novelty and exposure strategy.

\adds{} Using daily Hong Kong hospitalizations in hot seasons (May--October,
2000--2019), Guo et al.\ compared three nighttime metrics
\citep{guo2024hotnights}:
\begin{itemize}
  \item \textbf{HNday28$^\circ$C}: official hot night (Tmin $\ge$ 28$^\circ$C);
  \item \textbf{HNe}: hot-night excess (hourly degrees above 28$^\circ$C at night,
        $^\circ$C$\cdot$h);
  \item \textbf{HNday90th}: nights above the 90th percentile of HNe.
\end{itemize}
Extreme HNe was associated with higher non-cancer non-external hospitalizations
(about 3.1\% over lags 0--4), with larger effects in older and lower-SES groups and
for circulatory admissions. Critically, the official binary HNday28$^\circ$C metric
did \emph{not} show an overall excess-risk association, whereas intensity-based
metrics did.

\inspires{} Distinguish \emph{policy metrics} from \emph{biological metrics}.
Official binary hot nights remain policy-aligned and communicable. Intensity and
spell metrics may be closer to physiological stress. Consecutive intensive hot
nights mattered; late-season unusual nighttime heat was especially hazardous in
their interpretation.

\debates{} We cannot claim to be the first Hong Kong hot-night hospitalization
study. We also cannot treat a null finding for monthly official hot-night counts
as proof that nighttime heat is harmless. Guo's result predicts that binary counts
may be weak, while intensity/spell alternatives may be more informative.

\impl{}
\begin{enumerate}
  \item Keep official monthly hot-night counts as the \emph{primary policy-aligned}
        heat exposure.
  \item Prespecify intensity/spell alternatives as co-equal scientific checks, not
        post-hoc rescue analyses.
  \item Request hourly temperature (or precomputed HNe) from Hogan if feasible;
        daily Tmin alone cannot reconstruct true HNe.
  \item Until hourly data arrive, implement daily-derived spell summaries already
        foreshadowed in \texttt{scripts/02\_build\_weather\_monthly.R}
        (longest run; days in spells $\ge$3 or $\ge$5; 2D3N-style event flags).
  \item In manuscripts, state explicitly: ``A null official-hot-night association
        would be compatible with Guo et al.\ (2024) and would not close the
        nighttime-heat question.''
\end{enumerate}

\subsection{Wang et al.\ (2019) --- extremely hot weather events as combinations}

\adds{} Using official VHD (Tmax $\ge$ 33$^\circ$C) and HN (Tmin $\ge$ 28$^\circ$C)
definitions, Wang et al.\ showed that event \emph{structure} matters for mortality:
single hot nights raised risk; five or more consecutive VHDs/HNs raised it further;
combined patterns such as 2D3N (two very hot days with three hot nights) were
especially informative; females and older adults were more vulnerable
\citep{wang2019ehwe}.

\inspires{} Count-based monthly totals discard sequence. Two months with ten hot
nights can differ radically if one has scattered nights and the other has a
five-night siege. Spell and combination metrics recover some of that information
even after monthly aggregation.

\debates{} Mortality $\neq$ admissions. Warning-system behavioural responses may
also differ. Still, the exposure engineering lesson transfers cleanly.

\impl{} Extend monthly climate outputs with:
\begin{itemize}
  \item \texttt{days\_in\_hot\_night\_spell\_ge5}
  \item \texttt{days\_in\_very\_hot\_spell\_ge5}
  \item \texttt{longest\_cold\_run}
  \item a monthly indicator for whether a 2D3N-type event touches the month
\end{itemize}
These are already partially scaffolded; they should be first-class alternative
exposures in modelling scripts, not orphan columns.

%----------------------------------------------------------------------
\section{Pollution, influenza, and the cold pathway}
%----------------------------------------------------------------------

\subsection{Guo et al.\ (2025) --- temperature effects amplified on polluted days}

\adds{} Over 2000--2019, PM$_{2.5}$ amplified temperature--hospitalization
associations in Hong Kong; NO$_2$ intensified circulatory risks; risks were higher
on days with one or more elevated pollutants than on unpolluted days
\citep{guo2025temppollution}.

\inspires{} Pollution is not only a confounder checklist item. It can modify
temperature effects. Staged models should include interaction or stratified
``polluted month'' specifications once real EPIC series exist.

\debates{} Ozone remains ambiguous: correlated with hot sunny weather, possibly
pathway rather than pure confounder. Adjusting for O$_3$ in the primary heat model
may attenuate the very association of interest.

\impl{} Preserve staged pollution modelling:
(1) no pollutant;
(2) single-pollutant (NO$_2$, PM$_{2.5}$, O$_3$ separately);
(3) multi-pollutant sensitivity;
(4) optional temperature $\times$ pollution product terms.
Do not put O$_3$ into the default primary heat model without reporting the
unadjusted heat estimate beside it. Consider an Ox-style redox summary
($\frac{1}{3}$NO$_2 + \frac{2}{3}$O$_3$) as an exploratory combined oxidant metric
once real pollution data arrive \citep{wang2025coldflu}.

\subsection{Yang/Wei/Chong et al.\ (2025) --- cold, influenza, and stroke}

\adds{} In a 22-year weekly Hong Kong analysis ($\sim$1.17 million stroke-related
admissions), cold temperature and influenza activity (ILI+ rates) were both
associated with higher stroke admissions \citep{wang2025coldflu}. Absolute humidity
was preferred over relative humidity as a humidity metric less entangled with
temperature.

\inspires{} Influenza is a scientifically serious co-exposure for cold-season
stroke, not a decorative covariate. Absolute humidity is a better humidity
construct than RH for some infection/cold pathways.

\debates{} Weekly pooled stroke (ICD-9 430--434, primary and secondary diagnoses)
is not our monthly principal-diagnosis AMI/IS/HS design. Their finding strengthens
cold + flu hypotheses but does not license pooled-stroke primary analysis.

\impl{}
\begin{itemize}
  \item Keep influenza in sensitivity models initially (assumption A23), with a
        clear path to promote it if CHP series are strong.
  \item Build absolute humidity from temperature and RH (Tetens vapour pressure)
        in the weather pipeline; do not rely on RH alone.
  \item Request strain-aware ILI+ if available; otherwise a total ILI+ or
        consultation proxy is still better than nothing.
\end{itemize}

%----------------------------------------------------------------------
\section{Warm-climate winter excess and the ``heat replaced cold'' myth}
%----------------------------------------------------------------------

\subsection{Sipil{\"a} et al.\ (2017) and Lorking et al.\ (2020)}

\adds{} Stroke seasonality shows autumn/winter excess even in temperate Finland
\citep{sipila2017} and winter excess in tropical Thailand \citep{lorking2020}.
Warm climates are not heat-only climates for cerebrovascular burden.

\inspires{} Our framing should be coexistence, not succession: emerging heat
burden \emph{and} persistent cold risk in an aging subtropical city.

\debates{} Seasonality studies are not temperature-extreme studies. They support
motivation, not exposure definitions.

\impl{} Retain cold days as co-primary. In the research note and Introduction,
replace any ``heat replacing cold'' language with coexistence language. The
validated HKO figure already shows cold days persisting through high-heat years
(e.g., 13 cold days in 2021 alongside 61 hot nights).

%----------------------------------------------------------------------
\section{Outcome validity: why principal diagnosis is non-negotiable}
%----------------------------------------------------------------------

\subsection{Cho et al.\ (2024) and Liu et al.\ (2025)}

\adds{} East Asian administrative-data validation work shows that carefully defined
AMI and stroke algorithms can achieve high positive predictive values for first
events (around 90\% in the Korean NHIS validation), with lower but still useful
PPVs for recurrent events \citep{cho2024korea,liu2025tianjin}. Episode construction
matters.

\inspires{} Insist on principal diagnosis, episode/transfer rules, and first versus
recurrent distinctions where feasible.

\debates{} Hong Kong HA coding may be ICD-9-CM, ICD-10, or mixed across 2013--2023.
External PPVs do not automatically transfer.

\impl{} Data request to Roro should treat principal diagnosis, admission type,
transfer linkage, and coding system-by-year as hard requirements. Small-cell
suppression must never be coded as zero. If only secondary-diagnosis-inclusive
extracts are available, demote claims accordingly.

%----------------------------------------------------------------------
\section{Methodological scaffolding: what good time-series practice demands}
%----------------------------------------------------------------------

\subsection{Bhaskaran et al.\ (2013) and Gasparrini DLNMs}

\adds{} Time-series environmental epidemiology requires explicit handling of
seasonality, long-term trend, time-varying confounders, lagged associations, and
diagnostics \citep{bhaskaran2013}. DLNMs flexibly capture nonlinear
exposure--lag--response surfaces when daily data exist \citep{gasparrini2010}.
Globally, cold still accounts for a large share of temperature-attributable
mortality even as heat receives more public attention \citep{gasparrini2015}.
Morbidity associations are often weaker and more heterogeneous than mortality
associations \citep{ye2012}.

\inspires{} Our monthly design is a deliberate trade-off: we gain alignment with
official monthly extreme counts, demographic offsets, and hospital-burden
interpretation; we lose daily lag resolution. That trade-off should be stated as a
feature of the question, not apologized for as a defect.

\debates{} If HA ever releases daily aggregates, a companion daily DLNM would be
scientifically powerful---especially for testing whether Guo-style HNe effects
appear for principal-diagnosis AMI/IS/HS. Until then, do not overclaim lag biology.

\impl{} Current NB models with month fixed effects, smooth time trend, COVID
phases, and population $\times$ days offsets are consistent with Bhaskaran-style
discipline at monthly resolution. Diagnostics (residual ACF/PACF, overdispersion
checks) should remain mandatory before any real-data inference. Consider a future
optional daily module rather than forcing DLNM syntax onto 132 monthly points.

%----------------------------------------------------------------------
\section{Synthesis: what remains defensible for this project}
%----------------------------------------------------------------------

After reading the above as a system rather than a pile of citations, the
defensible contribution is:

\begin{quote}
Among Hong Kong residents aged 35+, estimate \textbf{monthly} associations of
\textbf{officially defined} thermal extremes (especially hot nights and cold days)
with \textbf{principal-diagnosis inpatient} AMI, ischemic stroke, and hemorrhagic
stroke during \textbf{2013--2023}, using age--sex population offsets, separate older
age bands, COVID-period handling, and staged pollution adjustment---without
assuming that heat has replaced cold, and without claiming primacy over recent
daily hot-night hospitalization research.
\end{quote}

\subsection{Prespecified hypotheses (recommended)}

\begin{enumerate}
  \item Cold days are positively associated with monthly AMI and ischemic-stroke
        admission rates.
  \item Official hot-night counts may show weak or null associations (compatible
        with Guo 2024); spell/intensity alternatives may be more informative,
        especially in older age groups.
  \item Associations differ by diagnosis (AMI vs IS vs HS).
  \item Ozone adjustment may attenuate heat associations if ozone partly lies on
        the pathway; total and pathway-adjusted estimates should both be reported.
  \item Influenza adjustment may attenuate cold--stroke associations in
        sensitivity models.
\end{enumerate}

\subsection{Claims to avoid}

\begin{itemize}
  \item ``First Hong Kong temperature--CVD study.''
  \item ``First Hong Kong hot-night hospitalization study.''
  \item ``Little is known about hot nights and cardiovascular admissions in Hong Kong.''
  \item ``We will demonstrate a regime shift from cold to heat.''
  \item Direct numeric comparison of monthly RRs with Goggins daily per-$^\circ$C
        estimates.
\end{itemize}

%----------------------------------------------------------------------
\section{From literature to pipeline: implementation backlog}
%----------------------------------------------------------------------

The following items translate the review into repository work. Items marked
\textbf{[done/partial]} already exist in some form; others are next steps.

\begin{center}
\renewcommand{\arraystretch}{1.15}
\begin{tabularx}{\textwidth}{@{}>{\raggedright\arraybackslash}p{3.2cm} X >{\raggedright\arraybackslash}p{2.4cm}@{}}
\toprule
\textbf{Literature cue} & \textbf{Pipeline / design action} & \textbf{Where} \\
\midrule
Guo 2024 intensity vs binary
  & Keep official hot nights primary; add spell/intensity alternatives; write null-interpretation rule
  & \texttt{config.yml}; Methods; script 11 \\
Wang 2019 5HN / 2D3N
  & Export monthly spell-ge5 and 2D3N-touch indicators as modelling covariates
  & script 02; variable dictionary \\
Yang/Chong 2025 absolute humidity
  & Compute monthly mean absolute humidity from $T$ and RH (Tetens)
  & \texttt{utils.R}; script 02 \\
Yang/Chong 2025 influenza
  & Load CHP ILI series; flu sensitivity model
  & \texttt{data\_raw/chp\_flu}; scripts 06, 11 \\
Guo 2025 pollution modification
  & Staged pollutant models + optional polluted-month interactions; cautious O$_3$
  & scripts 04, 11; assumption ledger A31 \\
Goggins 2012 subtypes
  & Mandatory diagnosis-stratified fits; no pooled stroke primary
  & script 11 \\
Lam 2018 diabetes
  & Request DM flag if feasible; else document limitation
  & Roro memo; ledger A27 \\
Cho/Liu coding validity
  & Principal diagnosis + episode rules as hard extract requirements
  & Roro memo \\
Bhaskaran diagnostics
  & Residual ACF/PACF and dispersion checks before inference
  & script 12 \\
Gasparrini daily DLNM
  & Optional future daily module if HA daily aggregates become available
  & new script later \\
\bottomrule
\end{tabularx}
\end{center}

\subsection{Creative but disciplined extensions}

Three ideas are worth considering without diluting the core paper:

\begin{enumerate}
  \item \textbf{Burden decomposition.} Even before causal claims, report how much of
        year-to-year variation in extreme counts is concentrated in multi-day
        spells. This is descriptive epidemiology that links HKO operations to
        hospital planning.
  \item \textbf{Care-seeking versus biology under COVID.} Use all-cause emergency
        admissions (if obtained) as a utilization tracer. Temperature associations
        that vanish only in 2020--2022 may reflect access disruption rather than
        biology.
  \item \textbf{Policy metric audit.} Treat the project partly as an audit of
        whether official extreme-day definitions, designed for warnings, also track
        diagnosis-specific hospital burden at monthly scale. That framing remains
        novel even after Guo 2024, because Guo's null for HNday28 was daily,
        hot-season, and not principal-diagnosis AMI/IS/HS with demographic offsets.
\end{enumerate}

%----------------------------------------------------------------------
\section{Search documentation and limitations of this review}
%----------------------------------------------------------------------

\begin{itemize}
  \item \textbf{Date:} 2026-07-10 (this pass).
  \item \textbf{Sources:} PubMed, PMC, publisher DOI pages (Lancet RH--WP, ACS ES\&T,
        Springer, PLOS, STOTEN), Epidemiology and Health, prior project evidence
        matrix.
  \item \textbf{Priority:} Hong Kong temperature--CVD; hot nights; stroke
        seasonality; coding validity; time-series methods.
  \item \textbf{Not claimed:} Exhaustive systematic review with dual screening and
        PRISMA flow. Living evidence matrix:
        \texttt{literature/evidence\_matrix.csv}.
  \item \textbf{Human follow-up:} complete author list for Liu et al.\ (2025) from
        PDF; confirm HA coding system before locking ICD algorithms.
\end{itemize}

\section*{Bottom line}

The literature does not ask us to abandon hot nights. It asks us to stop treating
the official binary count as self-justifying, to keep cold in the foreground, to
separate stroke subtypes, to stage pollution carefully, and to interpret monthly
results as hospital-burden evidence rather than daily-trigger evidence. Those
constraints make the project more honest---and, if executed well, more useful.

\bibliographystyle{plainnat}
\bibliography{../../literature/references}

\end{document}
